\scanpage{136}
\[
 \mathbb Z^2\longrightarrow H \tag*{(18)}
\]
compatible avec les orientations, i.e. tel que l'on ait
\[
 z_1\wedge z_2=\omega \tag*{(19)}
\]
où $z_1,z_2$ sont les images par (18) des éléments de la base canonique de $\mathbb Z^2$. Cet isomorphisme (18) définit un isomorphisme correspondant
\[
 \mathbb R^2\simeq V \tag*{(20)}
\]
d'où des isomorphismes
\[
 \left\{\begin{aligned}
 \mathbb P^1_{\mathbb R}&\simeq\mathbb P^1(V),\quad\mathbb P^1(\mathbb C)\simeq\Sigma_V\\
 \mathscr D&\simeq\Sigma_{V,\omega}^*
 \end{aligned}\right. \tag*{(21)}
\]
où $\mathscr D$ est le demi-plan de Poincaré, i.e. \deleted{l'ens} la variété complexe des $z\in\mathbb C$ tels que $\deleted{\Im z}\Im z>0$
\[
 \mathscr D=\{z\in\mathbb C\mid\Im z>0\}, \tag*{(22)}
\]
identifié à $\Sigma_{\mathbb R^2,\omega_0}^*$ en identifiant $z\in\mathscr D$ à la droite
\[
 L_z\deleted{\unsure{\mathbb P_z}}=\mathbb C\cdot\underbrace{(z,1)}_{ze_1+e_2}. \tag*{(23)}
\]
% slanted
\margin{$(e_1,e_2)$ base canonique de $\mathbb R^2$. Si on associait à $z\in\mathbb C$ la droite $\mathbb C(1,z)$, alors \unsure{à|ce serait au|correspond} $\Sigma^*_{\mathbb R^2,\omega_0}$ \unsure{correspondrait|correspondant|donnerait} le \deleted{\unsure{demi}} demi-plan $\mathscr D^-$ au lieu de $\mathscr D=\mathscr D^+$}
\margin{Ainsi, les classes d'isomorphie de courbes elliptiques avec \unsure{rigidif-|rigidité} de Torelli sont en corr. biunivoque avec l'ens des pts de $\mathscr D$.}

Sur l'ens des classes d'isomorphie des courbes elliptiques $E=(\Pi,\omega,L\subset H\otimes_{\deleted{\unsure{\mathbb R}}\mathbb Z}\mathbb C)$, avec rigidification de Torelli $\mathbb Z^2\simeq H$, le groupe $\mathrm{Sl}(2,\mathbb Z)$ opère à droite de façon évidente, par composition \inserted{à droite} de (18) avec les automorphismes de $\mathbb Z^2$. \deleted{\unsure{Quand on|Quant}}\inserted{On va}
% archive sheet number: 135
